\documentclass[nonewpage]{xourse}


% Start \part op een nieuwe pagina (Ximera doet dat niet per default)
\xpretocmd{\part}{\clearpage}{}{}

\newcommand{\richting}{Wetenschappen 6de jaar}%De klas/richting --> hoofding header
% \documentclass{xourse}
%in
%\title{\includegraphics[width=0.7\textwidth]{pic5/elektromagnetisme.jpg}\\[2cm]Fysica 5de jaar - Wetenschappen}
\title{Fysica 6de jaar - Wetenschappen}
\author{Daan Lipkens}
% titelblad (PDF): zie xmPreamble.tex, sectie Titelblad en QR-codes
\schooljaar{2026-2027}
% titelfoto: LIGO Hanford (zwaartekrachtgolven): trillingen, golven, gravitatie
\titelfoto{pic6/ligo_hanford.jpg}
\titelfotobron{NOIRLab/LIGO/NSF/AURA/T. Matsopoulos}{CC BY 4.0}{https://commons.wikimedia.org/wiki/File:LIGO_at_Hanford,_Washington_(LIGO_Pic_09-CC).jpg}
\titelbijschrift{LIGO meet zwaartekrachtgolven met laserlicht in twee loodrechte armen van 4 km lang. De spiegels hangen aan slingers die trillingen van de bodem wegfilteren. Trillingen, golven en gravitatie zijn leerstof van dit jaar.}
\cursuswebsite{https://xerxes.ximera.org/sui/6de/Cursus_DDO_6/6de/}
\begin{document}
\begin{abstract}
    Cursus Fysica 6de jaar Wetenschappen domeinoverschrijvende doorstroom.
\end{abstract}
\maketitle
%\xmtitle{Fysica 5de jaar - Wetenschappen}{}
% \part voor delen
%\chapterstyle voor hoofdstukken >> \activity{hoofdstuk ....}
%\sectionstyle voor delen in het hoofdstuk >> \activity{met onderdelen}
% Inleiding, informatie, ...


% INFO
% toepassingen onder \begin{example}
% Quizjes als \question met \begin{selectAll}
% voorbeeld oefeningen met align x &= berekening && (uitleg)
% komma's met {,} in math mode


 \inleidingstart % deel zonder nummer, hoofdstuk 0 (zie xmPreamble)
 \part{Inleiding}
 \inleidingstop
% \chapterstyle
% \activity{inleiding_stand_63.tex}
% \sectionstyle
% \activity{inleiding_praktische_richtlijnen_63.tex}
 \activity{inleiding_hoe_fysica_te_studeren_6.tex}
% \activity{inleiding_vectorkleuren.tex}


%--------positie, snelheid en versnelling----------------
\part{Positie, snelheid en versnelling}
\activity{x_v_a_oef_uitdaging.tex}

%--------Wetten van Newton----------------
\part{Kracht en beweging} %parts nog aanpassen
%COULOMBKRACHT of ELEKTRISCHE KRACHTWERKING
%picture path: pic5/elektriciteit/coulomb/

\activity{wetten_van_newton_oef_uitdaging.tex}

% %--------MAGNETISME----------------
% \part{Elektromagnetisme}
% \part{Lorentzkracht}
% \part{Inductiespanning}

% \newpage
% %--------KERNFYSICA----------------



%---------APPENDIX---------------

\part{Appendix} %nummering anders

\chapterstyle
\activity{vectoren_inleiding.tex}

\sectionstyle
\activity{vectoren_scalar.tex}
\activity{vectoren_component.tex}
\activity{vectoren_som.tex}
\activity{vectoren_verschil.tex}
\activity{vectoren_scalaire_vermenigvuldiging.tex}
\activity{vectoren_vectorieel_product.tex}


\chapterstyle
%Gedeeld met de 5de: deze bestanden staan enkel in ../5de (13/11/26).
%Werkt ../ niet bij het bouwen? Zet dan de oude regels terug (standaardvorm_*.tex in deze map).
\activity{../5de/appendix_standaardvorm_inleiding.tex}

\sectionstyle
\activity{../5de/appendix_standaardvorm_si.tex}
\activity{../5de/appendix_standaardvorm_omzetten_eenheden.tex}
\activity{../5de/appendix_standaardvorm_bc.tex}
\activity{../5de/appendix_standaardvorm_oefeningen.tex}

%Vector rekenen
    %vectoren
    %cosinus regel

%nuclide tabel
%standaard vorm
%logaritme en exponentiële functies
%---------Bibliografie
\part{Bibliografie}

\end{document}
